% -------------------------------------------------------
% ASSIGNED TO: \textcolor{green}{Person 2}
% -------------------------------------------------------
% INSTRUCTION: Keep existing gap analysis content from Section 3.7.
% ADD the red-marked critique paragraphs below — the supervisor
% requires genuine critique of cited papers, not only gap identification.
% Convert first-person to third person throughout.
% -------------------------------------------------------

The literature reviewed in Section 3 reveals the existence of a diagnostic gap at the intersection
of model robustness and explanation quality.

Section 3.1 established that XAI exists in two cultures. Blue XAI has dominated the field in terms of adoption, 
producing methods that justify model outputs to end users. Red XAI, as formalised by Biecek and Samek (2024), 
reorients XAI as an active interrogation tool, one that questions whether a model deserves to be explained at 
all before any explanation is surfaced. The RED XAI philosophy is well-articulated in theory; what remains is a 
complete pipeline that implements this philosophy from overfitting diagnosis through to measurable improvement 
in symbolic explanation quality. This project fills that gap.
\\\\
The GDSC dataset, while invaluable as a pharmacogenomic resource, exhibits the P >> N structure that reliably 
induces overfitting in unregularised deep models (Partin et al., 2023). Models must process tens of thousands 
of gene expression features measured across a comparatively small number of unique cancer cell lines, even 
though these are tested against hundreds of drugs to form a large dataset of combinations (Wu et al., 2025). 
Klauschen et al. (2024) and Partin et al. (2023) demonstrate that clinicians cannot verify whether a model's 
predictions are grounded in genuine oncogenic biology or learned shortcuts without an explicit auditing 
mechanism. The literature in cancer AI has produced increasingly powerful predictive models but has not 
systematically addressed the question of whether those models are ready to be explained 
(Wicker et al., 2022; Ye \& Farnia, 2025). The distinction matters clinically: a biologically incoherent rule 
extracted from an overfitted model may appear statistically valid while pointing clinicians toward irrelevant 
or even harmful treatment decisions (Pahde et al., 2023; Klauschen et al., 2024).
\\\\
Section 3.3 reviewed the theoretical relationship between loss landscape geometry and generalisation, 
establishing that models converging to sharp minima exhibit brittle, poorly generalising representations 
(Li et al., 2018; Keskar et al., 2016). SAM was reviewed as an intervention that explicitly seeks flat minima 
during training (Springer et al., 2024). What the literature has not established, however, is whether the 
geometric improvement that SAM produces in the loss landscape translates into measurable improvement in the 
quality of symbolic rules extracted from the model. The connection between landscape geometry and explanation 
quality has been theorised but not empirically demonstrated in the pharmacogenomic context. This project 
provides that empirical test.
\\\\
Section 3.4 reviewed the ANN literature relevant to this project, covering the DeepCCDS architecture and the 
Clever Hans problem as documented by Pahde et al. (2023). The R2R framework represents the closest 
methodological precedent: it uses XAI to identify model artefacts, localises them, and corrects model behaviour 
iteratively. However, R2R operates on image data, uses LRP-based attribution heatmaps as its explanation medium, 
and corrects models through fine-tuning with artefact masks. It does not address the tabular genomic domain, 
does not use symbolic rule extraction as its explanation modality, and does not investigate whether 
regularisation applied at the optimisation stage can prevent the Clever Hans problem from arising rather than 
repairing it after the fact. This project is therefore differentiated from R2R in domain, explanation type, 
intervention mechanism, and causal direction.
\\\\
Section 3.5 reviewed the evolutionary algorithm literature, specifically GAs and NSGA-II as applied to rule 
extraction. The GPRO\_RX system demonstrated that multi-objective genetic algorithms guided by PROMETHEE-assisted 
ranking can produce rule sets that outperform single-objective extraction methods on the combined criteria of 
fidelity and comprehensibility (Yedjour et al., 2024). However, GPRO\_RX and the broader EA rule extraction 
literature evaluate their methods on models with no deliberate control over the model's generalisation state, 
rule extraction is applied to whatever model happens to have been trained, without investigating whether the 
model's level of overfitting mediates rule quality. The comparative value of GA versus NSGA-II for rule 
extraction from models at different points on the generalisation spectrum remains unexplored. This project 
introduces that control by extracting rules from both a deliberately overfitted baseline and a SAM-regularised 
model, isolating the contribution of model health to rule quality independently of the extraction algorithm used.
\\\\
Section 3.6 reviewed the quality framework for symbolic explanations, fidelity, comprehensibility, and 
stability as defined by Arrieta et al. (2020) and operationalised through PROMETHEE II ranking by 
Yedjour et al. (2024). Adebayo et al. (2018) showed through sanity checks that popular XAI methods can 
produce visually convincing explanations even for models that have memorised noise. To the best of our 
knowledge, the literature does not establish a quantitative relationship between the size of the generalisation 
gap and the degradation of rule quality metrics. Whether a smaller generalisation gap is a statistically 
reliable predictor of higher fidelity and lower rule complexity, the core hypothesis of this project, 
has not been tested. Section 3.6.4 further noted that even statistically valid rules require biological 
verification through pathway enrichment analysis to confirm clinical relevance, introducing a layer of 
grounding that purely metric-based XAI evaluation cannot provide.
\\\\
Taken together, the literature has established four things independently: that overfitting produces unreliable 
explanations (Pahde et al., 2023; Adebayo et al., 2018); that loss landscape geometry is a measurable indicator 
of generalisation quality (Li et al., 2018); that SAM can improve landscape geometry (Foret et al., 2021); and 
that evolutionary algorithms can extract high-quality symbolic rules from neural networks when properly 
configured (Yedjour et al., 2024). The contribution of this project is not any single component, the 
ANN baseline, the SAM optimiser, the GA, the NSGA-II, and the PROMETHEE ranking are all individually 
established. The contribution is the controlled pipeline that connects them, the experimental structure that 
isolates model health as the independent variable, and the empirical evidence it generates about the 
relationship between generalisation and explainability in the pharmacogenomic domain.

% -------------------------------------------------------
% NEW CONTENT REQUIRED (NICOG Slide 4: "Critical analysis of previous work used")
% -------------------------------------------------------

\textcolor{red}{
% ADD: One sentence of genuine critique per key paper.
% The existing Section 3.7 identifies gaps (what has not been done).
% The supervisor requires critique of what the cited papers themselves do poorly.
% Suggested additions — integrate these into the relevant paragraphs above:
While Yedjour et al.\ \cite{Yedjour2024} demonstrated that PROMETHEE-assisted multi-objective
genetic algorithm rule extraction outperforms single-objective methods, their evaluation was
conducted exclusively on the UCI benchmark datasets and was not tested on high-dimensional
genomic data, where the curse of dimensionality and the presence of batch effects present
substantially different challenges to rule search spaces.
\\\\
Springer et al.\ \cite{Springer2024} established that SAM produces more balanced feature
representations than conventional regularisers, but their analysis was limited to image
classification benchmarks and does not provide evidence of this effect on tabular pharmacogenomic
data or on the downstream quality of symbolic rule extraction.
\\\\
Pahde et al.\ \cite{Pahde2023}, while providing a compelling audit-and-fix framework through
Reveal to Revise, operate exclusively in the image classification domain using LRP-based
attribution heatmaps. Their correction mechanism is applied post hoc through fine-tuning with
artefact masks, whereas this project intervenes at the optimisation stage to prevent Clever Hans
behaviour from arising rather than correcting it after the fact.
\\\\
Li et al.\ \cite{Li2018} provide strong theoretical and empirical evidence connecting flat minima to
improved generalisation, but do not investigate how loss landscape geometry affects the quality of
downstream symbolic explanations, leaving the connection between landscape smoothness and
explanation fidelity as an empirically untested claim that this project is designed to address.
\\\\
Adebayo et al.\ \cite{Adebayo2018} demonstrated through sanity checks that popular attribution
methods can produce convincing explanations for models that have memorised noise, but their
study did not evaluate symbolic rule-based explanation methods or propose a regularisation
intervention as a corrective mechanism, leaving the question of how to prevent rather than detect
this behaviour unresolved.
}

\subsubsection{Contributions}

This project makes the following concrete contributions to the literature.
The first is methodological. It introduces and evaluates a RED XAI pipeline that uses the generalisation gap as 
a diagnostic gatekeeper for symbolic explanation quality in cancer pharmacogenomics – a combination of domain, 
explanation type, and diagnostic mechanism that, to the best of our knowledge, has not been previously 
investigated. 
\\\\
The second is empirical. It provides a controlled experimental test of whether SAM regularisation, by smoothing 
the loss landscape, produces a statistically significant improvement in the fidelity and comprehensibility of 
evolutionary rule extraction on GDSC drug sensitivity data. Statistical rigour is ensured through Wilcoxon 
signed-rank testing across multiple experimental runs, following the precedent established by Yedjour et al. 
(2024), with the number of runs to be confirmed against the requirements of the adopted methodology in 
Section 4. \textcolor{red}{[PLACEHOLDER – confirm run count- justificatio needed]} 
\\\\
The third is evaluative. It provides a comparative assessment of two evolutionary rule extraction approaches – 
single-objective GA and multi-objective NSGA-II – applied to both a baseline overfitted model and a 
SAM-regularised model, across a consistent evaluation framework informed by Arrieta et al. (2020) and ranked 
through PROMETHEE II. This comparative structure generates evidence about which extraction method is more 
resilient to model fragility, a question that has not been addressed in the existing EA-XAI literature.
The fourth is diagnostic. The loss landscape visualisation produced for both the baseline and SAM-regularised 
model constitutes a novel proposed indicator of explainability readiness – a geometric measure of whether a 
model is positioned to yield reliable explanations before any extraction is attempted. This moves the XAI 
quality question upstream, from the explanation evaluation stage to the model training stage, consistent with 
the RED XAI philosophy of interrogating the model rather than justifying it.
